\cotefile{161-3}
\scanpage{4}
% L1
\uline{Produits de topos}

% L2–3
$\Pi(E,F)$ \unsure{$=$|$\simeq$} faisceaux sur $E$ à valeurs dans $F$ i.e. foncteurs $E^o\to F$ \deleted{\unsure{respectant}}commutant aux $\varprojlim$, i.e. qui ont un adjoint à gauche

% L4–5
$\simeq$ faisceaux sur $F$ à valeurs dans $E$, i.e foncteurs $F^o\unsure{\longrightarrow|-}E$ commutant aux $\varprojlim$, i.e. qui ont un adjoint à g.

% L6–7
$\simeq$ foncteurs $E^o\times F^o\longrightarrow\mathrm{Ens}$ qui commutent \uline{en chaque} \uline{variable} aux $\varprojlim$ quelc.

% L8-left
\margin{Si $E\cong\widetilde C$}

% L8–13
$\simeq$ faisceaux sur $C$ à valeurs dans $F$, i.e. foncteurs $\varphi\colon C^o\longrightarrow F$ tels que pour tt $R\hookrightarrow X$ couvrant dans $C$, on a $\varphi(x)\xrightarrow{\sim}\varprojlim_{\unsure{C/R|c/R}}\varphi(z)$, i.e. tels que $\forall\unsure{y|Y}\in F$, $x\mapsto\mathrm{Hom}(\varphi(x),y)$ sur $C$ est un faisceau, \inserted{(\unsure{ou autres|ou plutôt} \unsure{repr.|resp.|syst.} de flèches)} i.e. foncteurs $\varphi^o\colon C\to F^o$ qui sont ``\uline{continus}'' [i.e. se prolongent en foncteurs $\unsure{\overline C|\widetilde C}\longrightarrow F^o$ commutant aux $\varinjlim$ quelcon.]

% L14-left
\margin{Si $F\cong\widetilde D$}

% L14–16
$\simeq$ faisceaux sur $D$ à valeurs dans $E$, i.e. foncteurs $\psi\colon D^o\longrightarrow E$ tels que --, i.e. foncteurs ``continus'' $\psi^o\colon D\to E^o$

% L17–19
$\simeq$ foncteurs $C^o\times D^o\longrightarrow(\mathrm{Ens})$ qui sont des ``bifaisceaux'' i.e. tels que ce soient des faisceaux séparément en chaque variable

% L20
\deleted{\unsure{Soient} $u\colon E\longrightarrow E'$, $v\colon F\longrightarrow F'$ \unsure{d}}

% L21
\uline{Prop. \textcircled{1}} \quad $\Pi(E,F)$ est un topos

% L22
On peut supposer $E=\widetilde C$, $F=\widetilde D$, $C$ et $D$ petits sites stables par $\varprojlim$ finies.

% L23
\unsure{a)} Cas des topologies discrètes : $\pi(\widehat C,\widehat D)\cong\widehat{C\times D}$ : c'est un topos

% L24–25
b) Plus généralement, $\pi(\widehat C,\widetilde D)\simeq\underline{\mathrm{Hom}}\,(C^o,\widetilde D)$ est un topos, \deleted{\unsure{et} de \unsure{même}} $\pi(\widetilde C,\widehat D)\simeq\underline{\mathrm{Hom}}\,(D^o,\widetilde C)$.

% L26–29
\deleted{\unsure{Le|Et}} L'inclusion naturelle $\pi(\widehat C,\widetilde D)\hookrightarrow\Pi(\widehat C,\widehat D)$ \inserted{($\varphi\mapsto i_D\circ\varphi$, où $i_D\colon\widetilde D\longrightarrow\widehat D$ l'inclusion)} est le foncteur image directe pour un morphisme de topos, dont le foncteur image inverse \deleted{\unsure{est}} $\pi(\widehat C,\widehat D)\longrightarrow\pi(\widehat C,\widetilde D)$ est donné par $\varphi\longmapsto a_D\circ\varphi$, où $a_D\colon\widehat D\longrightarrow\widetilde D$ est l'adjoint à g. de $i_D$.

% L30–32
Donc $\pi(\widehat C,\widetilde D)$ est un sous-topos (str. plein) dans $\pi(\unsure{\widehat C|C},\widehat D)$ et pour la raison symétrique $\pi(\widetilde C,\widehat D)$ aussi. Mais $\pi(\widetilde C,\widetilde D)$ est l'intersection de ces deux, donc c'est un \inserted{sous-}topos de $\pi(\widehat C,\widehat D)$

% foot
% archive sheet number: 3
